\section{Formalism}
\label{sec:formal}

Deployment questions require probabilities that can be interpreted from
traces. The operational question is simple: after an agent has spent
several steps reasoning, calling tools, and observing outputs, what is
the probability that it reaches success before failure? We represent
that ``think-act-observe'' loop with a fitted absorbing discrete-time
Markov chain (DTMC), following classical absorbing-chain and
software-reliability models~\cite{kemenysnell,cheung1980}: intermediate
behavior becomes transient execution states, the two terminal outcomes
are task success~$\oplus$ and fatal failure~$\ominus$, and reliability
becomes a first-passage question---how soon a run reaches the success
absorber, and whether failure arrives first. This view also
gives benchmark metrics a common semantics, since
pass$@k$~\cite{chen2021codex}, pass$^k$~\cite{taub2024}, and the
step-budget reliability decay curve (RDC) in the sense of
\S\ref{sec:unify}~\cite{rdc2025} are all restrictions of one success
first-passage distribution. We write $\ST$ for the transient state set: the discrete
labels used for intermediate reasoning, tool-use, and observation
behavior before a run absorbs into $\oplus$ or $\ominus$.

%% [REDUCTION] fig:concept_translation is suppressed for the 12-page body:
%% Definition~\ref{def:amc} and Fig.~\ref{fig:pipeline} carry the model.
%% Retained in source for the artifact.
\iffalse
\begin{figure}[!t]
  \centering
  \begin{tikzpicture}[
      >=Stealth,
      scale=0.9, every node/.style={transform shape},
      concept/.style={rectangle, rounded corners, draw=blue!70, fill=blue!5, line width=1pt, align=center, font=\small\sffamily, minimum width=2.4cm, minimum height=0.8cm},
      success/.style={rectangle, rounded corners, draw=green!60!black, fill=green!5, line width=1pt, align=center, font=\small\sffamily, minimum width=2.4cm, minimum height=0.8cm},
      fail/.style={rectangle, rounded corners, draw=red!70, fill=red!5, line width=1pt, align=center, font=\small\sffamily, minimum width=2.4cm, minimum height=0.8cm},
      trans/.style={->, draw=gray!80, line width=0.8pt, font=\scriptsize}
  ]
      % Transient nodes
      \node[concept] (plan) at (0, 0) {Reason / Plan};
      \node[concept] (tool) at (4.5, 0) {Call Tool API};
      \node[concept] (obs) at (2.25, -2) {Observe Output};
      
      % Absorbing nodes
      \node[fail] (fail1) at (0, -4.2) {Fatal Error ($\ominus$)};
      \node[success] (succ) at (4.5, -4.2) {Task Solved ($\oplus$)};
      
      % Transient edges
      \draw[trans] (plan) to[bend left=15] node[above] {execute} (tool);
      \draw[trans] (tool) to[bend left=15] node[left, near start] {result} (obs);
      \draw[trans] (obs) to[bend left=15] node[right, pos=0.55, xshift=2pt] {revise} (plan);
      
      % Absorbing edges
      \draw[trans, color=red!80] (plan) -- node[left] {crash} (fail1);
      \draw[trans, color=red!80] (tool.south) .. controls +(0.9,-1.8) and +(1.1,0.6) .. node[right, pos=0.72] {timeout} (fail1.east);
      \draw[trans, color=red!80] (obs) -- node[left, near start] {wrong} (fail1);
      \draw[trans, color=green!60!black] (obs) -- node[right] {correct} (succ);
      
      % Background box
      \begin{scope}[on background layer]
          \node[draw=blue!20, fill=blue!2, rounded corners=5pt, fit=(plan)(tool)(obs), inner sep=10pt] (tbox) {};
      \end{scope}
      \node[text=blue!60!black, font=\scriptsize\itshape, above=2pt of tbox] {Transient States Loop ($\ST$)};
      
  \end{tikzpicture}
  \caption{Conceptual translation from an agent's interactive loop to
           a DTMC. Reasoning phases become transient states, and each
           run eventually absorbs into success or failure.}
  \label{fig:concept_translation}
\end{figure}
\fi

Later sections estimate the components of $\hat M$ from traces and test
whether the abstraction is adequate before using it for reliability
claims.

\begin{definition}[Agent Markov chain]\label{def:amc}
An \emph{agent Markov chain} (AMC) is a tuple
$\hat M=(\ST,\hat Q,\hat R_\oplus,\hat R_\ominus,s_0)$, where $\ST$ is
the finite set of transient \emph{execution states} obtained from trace
featurization (\S\ref{sec:state_construction}),
$\hat Q\in[0,1]^{|\ST|\times|\ST|}$ gives transitions among those
states, and $\hat R_\oplus,\hat R_\ominus\in[0,1]^{|\ST|}$ give
absorption probabilities into the task-complete state $\oplus$ and the
fatal-failure state $\ominus$. Mass is conserved:
$\hat Q\mathbf{1}+\hat R_\oplus+\hat R_\ominus=\mathbf{1}$. Finally,
$s_0\in\ST$ is the initial state, or more generally an initial
distribution~$\pi_0$.
\end{definition}

\begin{definition}[Reliability]\label{def:reliab}
The $d$-step reliability is
$\mathcal R(d;\hat M)=\Pr[\tau_\oplus\le d\mid X_0=s_0]$, where
$\tau_\oplus$ is the first time the chain reaches the success
absorber, and $R_\infty=\lim_{d\to\infty}\mathcal R(d;\hat M)$. The
Kemeny--Snell fundamental matrix~\cite{kemenysnell} is
$N=(I-\hat Q)^{-1}$; it counts expected visits to transient states
before absorption and yields the closed form in
Prop.~\ref{thm:closed}.
\end{definition}

\noindent\textbf{Benchmark metrics.}
For $k$ repeated runs of a task, $\mathrm{pass}@k$ is the probability that
\emph{at least one} run succeeds and $\mathrm{pass}^k$ that \emph{all} $k$
succeed; under the fitted chain and i.i.d.\ trials (A\ref{A:iid}) both are
restrictions of the same success first-passage distribution as the RDC
(Theorem~\ref{thm:unify}).

\begin{table}[!t]
\centering
\caption{Notation. The fitted chain $\hat M$ and its diagnostics are
estimated from traces (\S\ref{sec:state_construction}).}
\label{tab:notation}
\small
\setlength{\tabcolsep}{4pt}
\begin{tabular}{@{}ll@{}}
\toprule
\textbf{Symbol} & \textbf{Meaning} \\
\midrule
$\ST$, $m$, $\phi$ & transient states, $m=|\ST|$, step featurizer \\
$\hat Q$ & $m{\times}m$ transient transition matrix \\
$\hat R_\oplus,\hat R_\ominus$ & exit vectors to success $\oplus$ / failure $\ominus$ \\
$N{=}(I{-}\hat Q)^{-1}$ & fundamental matrix (visits pre-absorption) \\
$s_0$, $\tau_\oplus$ & start state; first-passage time to success \\
$\mathcal R(d),R_\infty$ & reliability $\Pr[\tau_\oplus\!\le\! d]$ and its limit \\
$\rho(\hat Q)$ & spectral radius of $\hat Q$ (A\ref{A:trans}) \\
$\mathrm{pass}@k,\mathrm{pass}^k$ & any-of-$k$ / all-of-$k$ success probabilities \\
$\Delta_{\mathrm{AIC}}$ & AIC(1st)$-$AIC(2nd) Markov-order statistic \\
\bottomrule
\end{tabular}
\end{table}

This first-passage view also connects the fitted chain to classical
reliability growth. A central software-reliability model is the
non-homogeneous Poisson process (NHPP), a counting-process model for
failure discovery over time, especially the Goel--Okumoto
model~\cite{goelokt}; Musa--Okumoto gives a related formulation~\cite{musa1987}.
Under rare-failure scaling, the fitted agent chain recovers this NHPP
form as a limit (Proposition~\ref{thm:nhpp}), which explains when
classical reliability-growth estimators are meaningful for agent traces.

\noindent\textbf{Model assumptions.}
Transience gives eventual absorption and the fundamental matrix;
invertibility is the algebraic condition for the closed forms; the
initial condition fixes where traces start; and i.i.d.\ replications are
required when interpreting pass$@k$ and pass$^k$. The first-order Markov
property is not taken on faith: the tests in
\S\ref{sec:order_selection} and \S\ref{sec:gof} report when the corpus
does not support the fitted DTMC.

\begin{assumption}[Transience]\label{A:trans}
$\rho(\hat Q)<1$.
\end{assumption}
\begin{assumption}[Invertibility]\label{A:inv}
$(I-\hat Q)$ is invertible.
\end{assumption}
\begin{assumption}[Initial condition]\label{A:init}
$s_0$ is deterministic (extensions replace $e_{s_0}^\top$ by
$\pi_0^\top$).
\end{assumption}
\begin{assumption}[I.i.d.\ replications]\label{A:iid}
For pass$^k$ and pass$@k$, trials are i.i.d.\ unless noted.
\end{assumption}

\begin{definition}[Local perturbation family]\label{def:perturb}
For sensitivity analysis, let
$Q(\varepsilon)=Q_0+\varepsilon\Delta$ for
$0\le\varepsilon\le\varepsilon_{\max}$, where rows remain
substochastic and $\rho(Q(\varepsilon))<1$. The success-exit vector
$R_{\oplus,0}$ is held fixed, and any removed transient mass is
assigned to the failure absorber so that
$Q(\varepsilon)\mathbf{1}+R_{\oplus,0}+R_{\ominus}(\varepsilon)
=\mathbf{1}$. Operationally, this represents a local design change,
such as adding a fallback tool, without re-running the full benchmark.
\end{definition}

\noindent\textbf{Remark.}
The fitted absorbing DTMC is not intended to recover every hidden state
of the agent or its environment. It is a compact reliability model
estimated from observable traces, and its outputs are meaningful only
when the state construction and goodness-of-fit checks support the
first-order absorbing-DTMC approximation.

%% ----------------------------------------------------------------
%% [REDUCTION] The toy fitted-DTMC illustration (fig:markov_chain)
%% is suppressed from the 12-page body; Fig.~\ref{fig:concept_translation}
%% already shows the transient loop and the two absorbers.
%% ----------------------------------------------------------------
\iffalse
\begin{figure}[!t]
  \centering
  \begin{tikzpicture}[
      >=Stealth,
      scale=0.85, every node/.style={transform shape},
      transient/.style={circle, draw=blue!80, fill=blue!10, line width=1pt, minimum size=0.9cm, font=\small\bfseries},
      absorbing/.style={circle, draw=#1, fill=#1, fill opacity=0.15, text opacity=1, line width=1.5pt, minimum size=1cm, font=\small\bfseries, double, double distance=1.5pt, double=white},
      trans/.style={->, draw=gray!80, line width=0.8pt, font=\scriptsize},
      tosucc/.style={->, draw=green!60!black, line width=1pt, font=\scriptsize},
      tofail/.style={->, draw=red!70, line width=1pt, font=\scriptsize}
  ]
      % States
      \node[transient] (s0) at (0, 2.2)  {$s_0$};
      \node[transient] (s1) at (0, 0)    {$s_1$};
      \node[transient] (s2) at (0,-2.2)  {$s_2$};
      \node[absorbing=green!60!black] (succ) at (4.5, 1.1)  {$\oplus$};
      \node[absorbing=red!70]   (fail) at (4.5,-1.1)  {$\ominus$};

      % Transitions (Q)
      \draw[trans] (s0) edge[loop left, looseness=4, out=160, in=200] node[left] {0.5} (s0);
      \draw[trans] (s0) -- node[right] {0.3} (s1);
      \draw[trans] (s1) edge[loop left, looseness=4, out=160, in=200] node[left] {0.4} (s1);
      \draw[trans] (s1) -- node[right] {0.3} (s2);
      \draw[trans] (s2) edge[loop left, looseness=4, out=160, in=200] node[left] {0.6} (s2);

      % Transitions (R)
      \draw[tosucc] (s0) to[bend left=10]  node[above, sloped] {0.1} (succ);
      \draw[tofail] (s0) to[bend right=10] node[below, sloped] {0.1} (fail);
      \draw[tosucc] (s1) to[bend left=5]   node[above, sloped] {0.2} (succ);
      \draw[tofail] (s1) to[bend right=5]  node[below, sloped] {0.1} (fail);
      \draw[tosucc] (s2) to[bend left=10]  node[above, sloped] {0.3} (succ);
      \draw[tofail] (s2) to[bend right=10] node[below, sloped] {0.1} (fail);

      % Labels
      \begin{scope}[on background layer]
          \node[draw=blue!20, fill=blue!5, rounded corners=5pt, fit=(s0)(s1)(s2), inner sep=10pt] (tbox) {};
      \end{scope}
      \node[text=blue!60!black, font=\scriptsize\itshape, above=1pt of tbox] {Transient states $\mathcal{S}_T$};
  \end{tikzpicture}
  \caption{Illustration of a fitted absorbing DTMC. Transitions among
           transient states (matrix $Q$) model agent logic, while
           success and failure exit vectors $R_\oplus$ and
           $R_\ominus$ lead to absorbing states $\oplus$ and
           $\ominus$.}
  \label{fig:markov_chain}
\end{figure}
\fi
%% ================================================================
